\documentclass[12pt]{article}
\usepackage[utf8]{inputenc}
\usepackage[spanish]{babel}
\usepackage{amsmath}
\usepackage{amssymb}
\usepackage{geometry}
\geometry{a4paper, margin=1in}
\begin{document}
\section{Módulo: src/render (Representación Visual Exclusiva)}

El directorio \texttt{src/render} es el final de la cadena de montaje computacional (Pipeline). Es el módulo consumidor puro: no realiza cálculos físicos ni análisis de audio; su única misión es recopilar los tensores numéricos del \texttt{core}, asignarles texturas y enviarlos a la pantalla o a un archivo de video final de manera ultarrápida.

\subsection{Estructura del Motor de Renderizado}
El directorio se divide en ramas de estabilidad:
\begin{itemize}
    \item \textbf{\texttt{stable/render\_standard.py}:} El bucle principal de renderizado usando OpenCV (\texttt{cv2}). Se encarga de transformar las matrices abstractas del \texttt{core} en fotogramas RGB (\textit{frames}).
    \item \textbf{\texttt{experimental/}:} Contiene motores de renderizado más avanzados o no lineales (como trazado de rayos básico o Shaders nativos) que están en fase beta.
\end{itemize}

\subsection{Proceso de Dibujado (Drawing Pipeline)}
El motor \texttt{render\_standard.py} orquesta la función \texttt{generar\_animacion\_god\_mode()}. Su ciclo de vida es el siguiente:
\begin{enumerate}
    \item Lee la matriz escalar generada por \texttt{src/core} (por ejemplo, el mapa de densidad de un fluido).
    \item Aplica un mapeo de color (Colormap) dinámico basado en las variables dictadas por \texttt{src/audio} (Tonalidad).
    \item Si corresponde, dibuja las partículas usando operaciones de trazado matricial o \texttt{cv2.circle}.
    \item Si el \texttt{motor\_lyrics.py} envía texto sincronizado, invoca rutinas tipográficas (\texttt{Pillow} o \texttt{cv2.putText}) para sobreimprimir la lírica en pantalla con colores vibrantes y desenfoques (Bloom).
    \item Escribe el fotograma directamente en el \textit{VideoWriter} (FFmpeg) para exportar el `.mp4` final, reportando el progreso a la interfaz de usuario.
\end{enumerate}

\end{document}
